\documentclass[12pt]{article}
\usepackage[utf8]{inputenc}
\usepackage[T1]{fontenc}
\usepackage[margin=1in]{geometry}
\usepackage{setspace}
\usepackage{xcolor}
\usepackage[colorlinks=true, urlcolor=blue!60!black, linkcolor=blue!60!black]{hyperref}
\setlength{\parindent}{2em}
\setlength{\parskip}{1.25em}
\renewcommand{\baselinestretch}{1.0}
\begin{document}
\begin{center}
{\LARGE \textit{Seila Law LLC v. Consumer Financial Protection Bureau}, 591 U.S. 197 (2020)\par}
\vspace{0.8em}
{\large Case Brief\par}
\end{center}
\vspace{1em}

\section{Memory Jogger}

Under Article II's Vesting Clause, Congress violates the separation of powers when it vests significant executive authority in an independent agency headed by a single Director removable only for cause, though the offending removal restriction is severable from the remainder of the Dodd-Frank Act (\href{https://www.courtlistener.com/opinion/4764451/seila-law-llc-v-consumer-financial-protection-bureau/}{CourtListener}; \href{https://www.law.cornell.edu/supremecourt/text/19-7}{Cornell LII}).

\section{Detailed Case Facts}

In the aftermath of the 2008 financial crisis---which followed the collapse of the subprime mortgage market and wiped out over \$10 trillion in American household wealth---Congress enacted the Dodd-Frank Wall Street Reform and Consumer Protection Act of 2010 (124 Stat. 1376), establishing the Consumer Financial Protection Bureau (CFPB) as an independent regulatory bureau within the Federal Reserve System (\href{https://www.law.cornell.edu/supremecourt/text/19-7}{Cornell LII}). Drawing on a 2007 proposal by then-Professor Elizabeth Warren and subsequent recommendations from the Obama Treasury Department, Congress consolidated consumer financial oversight into a single agency tasked with ensuring that markets for consumer financial products and services are ``fair, transparent, and competitive'' (12 U.S.C. \S 5511(a)) (\href{https://www.law.cornell.edu/supremecourt/text/19-7}{Cornell LII}).

Congress vested the CFPB with sweeping rulemaking, enforcement, and adjudicatory authority. It transferred administration of 18 existing federal consumer-protection statutes to the Bureau---including the Fair Credit Reporting Act, the Fair Debt Collection Practices Act, and the Truth in Lending Act---and added a new statutory prohibition against ``any unfair, deceptive, or abusive act or practice'' in the consumer-finance sector (12 U.S.C. \S 5536(a)(1)(B)) (\href{https://www.law.cornell.edu/supremecourt/text/19-7}{Cornell LII}). The Bureau may promulgate binding regulations, conduct investigations, issue subpoenas and civil investigative demands (CIDs), initiate administrative adjudications, and prosecute civil enforcement actions in federal court seeking restitution, disgorgement, injunctive relief, and civil penalties of up to \$1,000,000 per day (inflation adjusted) (12 U.S.C. \S\S 5562--5565). By the time the case reached the Supreme Court, the CFPB had secured over \$11 billion in relief for more than 25 million consumers (\href{https://www.law.cornell.edu/supremecourt/text/19-7}{Cornell LII}).

Crucially, Congress departed from both Professor Warren's original proposal and the traditional model of multimember independent commissions. Instead of a bipartisan board, Congress placed the CFPB under a single Director appointed by the President with the advice and consent of the Senate for a five-year term (12 U.S.C. \S 5491(b)--(c)), during which the President could remove the Director only for ``inefficiency, neglect of duty, or malfeasance in office'' (12 U.S.C. \S 5491(c)(3)) (\href{https://www.law.cornell.edu/supremecourt/text/19-7}{Cornell LII}). Moreover, rather than relying on annual congressional appropriations, the CFPB draws its funding directly from the Federal Reserve---up to 12\% of the Federal Reserve's 2009 operating expenses, adjusted for inflation (exceeding \$500 million annually) (12 U.S.C. \S 5497(a)) (\href{https://www.law.cornell.edu/supremecourt/text/19-7}{Cornell LII}).

In 2017, during the tenure of first CFPB Director Richard Cordray, the Bureau issued a civil investigative demand under 12 U.S.C. \S 5562(c)(1) to Seila Law LLC, a California-based law firm providing debt-relief legal services, to investigate whether the firm had engaged in unlawful practices in advertising, marketing, or selling debt-relief services (\href{https://www.law.cornell.edu/supremecourt/text/19-7}{Cornell LII}). The CID directed Seila Law to produce documents and respond to interrogatories. Seila Law petitioned the CFPB to set aside the demand, asserting that the Bureau's leadership by a single Director removable only for cause violated the Constitution's separation of powers. The CFPB declined to address the constitutional objection and ordered compliance (\href{https://www.law.cornell.edu/supremecourt/text/19-7}{Cornell LII}).

\section{Procedural History}

When Seila Law refused to comply with the CID, the CFPB petitioned the U.S. District Court for the Central District of California in June 2017 to enforce the demand under 12 U.S.C. \S 5562(e)(1). On August 25, 2017, Judge Josephine L. Staton rejected Seila Law's separation-of-powers defense and ordered compliance with the CID (with one minor modification), holding that the CFPB's structure was constitutional (\textit{CFPB v. Seila Law LLC}, 2017 WL 6536586 (C.D. Cal. Aug. 25, 2017)) (\href{https://www.law.cornell.edu/supremecourt/text/19-7}{Cornell LII}; \href{https://fedsoc.org/case/seila-law-llc-v-consumer-financial-protection-bureau}{Federalist Society}).

On appeal, the U.S. Court of Appeals for the Ninth Circuit affirmed (\textit{CFPB v. Seila Law LLC}, 923 F.3d 680 (9th Cir. 2019)) (\href{https://www.law.cornell.edu/supremecourt/text/19-7}{Cornell LII}). Expressing agreement with the en banc D.C. Circuit's ruling in \textit{PHH Corp. v. CFPB}, 881 F.3d 75 (D.C. Cir. 2018), the Ninth Circuit acknowledged that the CFPB wielded more executive power than the 1935 FTC and was headed by a single Director rather than a board, but concluded that \textit{Humphrey's Executor v. United States}, 295 U.S. 602 (1935), and \textit{Morrison v. Olson}, 487 U.S. 654 (1988), were controlling and foreclosed Seila Law's challenge (\href{https://www.law.cornell.edu/supremecourt/text/19-7}{Cornell LII}).

Seila Law petitioned for a writ of certiorari. Notably, by the time the case reached the Supreme Court, both the Solicitor General and newly appointed CFPB Director Kathleen Kraninger agreed with Seila Law on the merits that the for-cause removal restriction in 12 U.S.C. \S 5491(c)(3) was unconstitutional, while maintaining that the provision was severable and that the CID remained enforceable because Acting Director Mick Mulvaney (who was removable at will) had ratified it (\href{https://www.law.cornell.edu/supremecourt/text/19-7}{Cornell LII}). The Supreme Court granted certiorari in October 2019 on both the separation-of-powers question and the severability question, and appointed former Solicitor General Paul Clement as \textit{amicus curiae} to defend the judgment below. The case was argued on March 3, 2020, and decided on June 29, 2020 (\href{https://www.law.cornell.edu/supremecourt/text/19-7}{Cornell LII}; \href{https://supreme.justia.com/cases/federal/us/591/19-7/}{Justia}).

\section{Judicial Votes}

The Court split differently on the constitutional merits (5--4) and on severability (7--2) (\href{https://www.law.cornell.edu/supremecourt/text/19-7}{Cornell LII}; \href{https://www.supremecourt.gov/opinions/19pdf/19-7_n6io.pdf}{Supreme Court Slip Opinion}):

\begin{itemize}
\item \textbf{Removal / Separation-of-Powers Holding (Parts I, II, and III) --- 5--4:}
\begin{itemize}
\item \textbf{Majority (5):} Chief Justice Roberts (author), joined by Justices Thomas, Alito, Gorsuch, and Kavanaugh.
\item \textbf{Dissenting on Removal (4):} Justice Kagan (author), joined by Justices Ginsburg, Breyer, and Sotomayor.
\end{itemize}
\item \textbf{Severability Holding (Part IV and Judgment) --- 7--2:}
\begin{itemize}
\item \textbf{Plurality (3):} Chief Justice Roberts (author), joined by Justices Alito and Kavanaugh.
\item \textbf{Concurring in the Judgment as to Severability (4):} Justice Kagan (author), joined by Justices Ginsburg, Breyer, and Sotomayor (agreeing that, given the majority's constitutional holding, the removal restriction must be severed rather than invalidating the entire CFPB).
\item \textbf{Dissenting on Severability (2):} Justice Thomas (author), joined by Justice Gorsuch (declining to reach severability and voting simply to deny enforcement of the CID).
\end{itemize}
\end{itemize}

\section{Holding}

\textbf{Judgment:} The judgment of the U.S. Court of Appeals for the Ninth Circuit is vacated, and the case is remanded for further proceedings to determine whether the civil investigative demand issued to Seila Law was validly ratified by an Acting Director accountable to the President (\href{https://www.law.cornell.edu/supremecourt/text/19-7}{Cornell LII}).

\textbf{Rule Announced:} Article II vests the entire ``executive Power'' in the President, which as a general rule includes the unrestricted power to supervise and remove executive officers at will. Existing precedent recognizes only two narrow exceptions: (1) multimember expert agencies balanced along partisan lines that do not wield substantial executive power (\textit{Humphrey's Executor v. United States}, 295 U.S. 602 (1935)), and (2) inferior officers with limited duties and no policymaking or significant administrative authority (\textit{Morrison v. Olson}, 487 U.S. 654 (1988)). Congress violates the separation of powers by extending for-cause removal protection to a single principal officer who heads an independent agency wielding significant executive power. However, under settled severability principles and Dodd-Frank's express severability clause (12 U.S.C. \S 5302), the unconstitutional removal restriction in 12 U.S.C. \S 5491(c)(3) is severable from the rest of the Act, leaving the CFPB operational with a Director removable by the President at will (\href{https://www.law.cornell.edu/supremecourt/text/19-7}{Cornell LII}).

\textbf{Holding vs. Dicta:} The binding \textit{holding} invalidates the for-cause removal protection for the single CFPB Director (5--4) and severs \S 5491(c)(3) from the Dodd-Frank Act so the Bureau may continue operating (7--2), while rejecting threshold standing, ripeness, and adverseness objections. Several important passages are \textit{dicta} or expressly reserved: (1) the Court expressly declined to revisit or overrule \textit{Humphrey's Executor} or \textit{Morrison}; (2) the Court's discussion of the CFPB's non-appropriated funding through the Federal Reserve and five-year term was explicitly labeled an aggravating feature rather than an independent ground of unconstitutionality (the funding mechanism was later upheld against an Appropriations Clause challenge in \textit{CFPB v. Community Financial Services Ass'n of America}, 601 U.S. 416 (2024)) (\href{https://supreme.justia.com/cases/federal/us/601/22-448/}{Justia}); (3) the Court's skepticism toward other single-head independent agencies---specifically the Office of Special Counsel, the Social Security Administration, and the Federal Housing Finance Agency---was dicta as to those entities (though the Court soon applied \textit{Seila Law} to invalidate the FHFA Director's removal protection in \textit{Collins v. Yellen}, 594 U.S. 220 (2021)) (\href{https://www.law.cornell.edu/wex/unitary_executive_theory_\%28uet\%29}{Cornell LII Wex}); and (4) the Court declined to decide whether Acting Director Mulvaney's ratification cured the constitutional defect in the CID.

\section{Analysis of Opinions}

\subsection{Opinion of the Court (Parts I--III) and Plurality Opinion (Part IV) --- Chief Justice Roberts}

\textbf{Threshold Justiciability (Part II):} Chief Justice Roberts rejected three threshold objections raised by court-appointed \textit{amicus} Paul Clement (\href{https://www.law.cornell.edu/supremecourt/text/19-7}{Cornell LII}). First, Seila Law had Article III appellate standing because it suffered a concrete injury from being compelled to comply with the CID; under \textit{Free Enterprise Fund v. PCAOB}, 561 U.S. 477 (2010), and \textit{Bowsher v. Synar}, 478 U.S. 714 (1986), a challenger need not prove that the agency would have acted differently in a ``counterfactual world'' without the removal restriction. Second, a removal restriction need not await a contested presidential firing to be ripe, as it inflicts a ``here-and-now'' injury on regulated parties. Third, despite the Government's agreement on the constitutional merits, the case retained Article III adverseness under \textit{United States v. Windsor}, 570 U.S. 744 (2013), because the Government still sought to enforce the CID on remand, and appointing \textit{amicus} satisfied prudential concerns (\href{https://www.law.cornell.edu/supremecourt/text/19-7}{Cornell LII}).

\textbf{Constitutional Merits (Part III):} Anchoring his analysis in Article II's Vesting Clause (``The executive Power shall be vested in a President'') and Take Care Clause, the Decision of 1789 in the First Congress, \textit{Myers v. United States}, 272 U.S. 52 (1926), and \textit{Free Enterprise Fund}, Roberts framed unrestricted presidential removal power as the constitutional rule, not the exception (\href{https://www.law.cornell.edu/supremecourt/text/19-7}{Cornell LII}). He cabined the Court's prior removal cases into two narrow exceptions: \textit{Humphrey's Executor} (multimember expert boards balanced along partisan lines with staggered terms exercising ``quasi-legislative'' or ``quasi-judicial'' functions, and said in 1935 to exercise no executive power) and \textit{Morrison} / \textit{United States v. Perkins}, 116 U.S. 483 (1886) (inferior officers with limited duties and no policymaking or administrative authority).

Roberts concluded that the CFPB fit neither exception. Unlike the 1935 FTC, the CFPB is led by a single Director who is neither a bipartisan ``body of experts'' nor a mere legislative or judicial aid, and who exercises quintessential executive power by promulgating binding rules across 19 statutes, issuing final adjudicatory orders, and seeking massive civil penalties in federal court. Unlike the independent counsel in \textit{Morrison}, the CFPB Director is a principal officer with vast policymaking and enforcement power aimed outward at millions of private citizens and businesses (\href{https://www.law.cornell.edu/supremecourt/text/19-7}{Cornell LII}).

Declining to extend those exceptions to a ``new situation,'' Roberts emphasized two fatal flaws. First, the CFPB's single-Director independent structure lacked historical footing: the four cited predecessors (the Civil War-era Comptroller of the Currency for a single year, the Office of Special Counsel, the Social Security Administration, and the FHFA) were either brief aberrations or modern, contested, and far narrower in private enforcement authority. Second, the structure clashed with constitutional design: the Framers divided power everywhere except the Presidency, fortifying a single President with energy and dispatch while making him uniquely accountable to the entire electorate. Vesting unilateral executive power in a single unaccountable Director---aggravated by a five-year term that can prevent a President from ever appointing a Director, and by funding outside the annual appropriations process---violated that chain of dependence. Roberts also refused to invoke constitutional avoidance to read ``inefficiency, neglect of duty, or malfeasance in office'' as permitting policy-based removal, noting that \textit{Humphrey's Executor} had rejected that reading and Dodd-Frank expressly styled the CFPB an ``independent bureau'' (\href{https://www.law.cornell.edu/supremecourt/text/19-7}{Cornell LII}).

\textbf{Severability (Part IV, Plurality):} Joined by Justices Alito and Kavanaugh, Roberts concluded that \S 5491(c)(3) was severable from Dodd-Frank. Relying on \textit{Alaska Airlines, Inc. v. Brock}, 480 U.S. 678 (1987), \textit{Free Enterprise Fund}, and Dodd-Frank's express severability clause (12 U.S.C. \S 5302), Roberts reasoned that the CFPB's remaining provisions remained fully operative and that eliminating the entire 1,500-employee agency would cause massive regulatory disruption (especially since the Office of Thrift Supervision had been abolished). Courts must ``use a scalpel rather than a bulldozer,'' leaving Congress free to redesign the CFPB as a multimember commission if it wishes (\href{https://www.law.cornell.edu/supremecourt/text/19-7}{Cornell LII}).

\subsection{Concurrence in Part and Dissent in Part --- Justice Thomas (Joined by Justice Gorsuch)}

Justice Thomas joined Parts I--III because the Court rightly cabined \textit{Humphrey's Executor} to ``multimember expert agencies that do not wield substantial executive power,'' but urged the Court to overrule \textit{Humphrey's Executor} altogether in a future case (\href{https://www.law.cornell.edu/supremecourt/text/19-7}{Cornell LII}). In Thomas's view, \textit{Humphrey's Executor} rested on the constitutionally impossible fiction of ``quasi-legislative'' and ``quasi-judicial'' powers floating outside the three vested branches of Articles I, II, and III, creating an unaccountable ``de facto fourth branch'' that threatens individual liberty. Because modern precedent (\textit{Morrison}, \textit{City of Arlington v. FCC}, and \textit{Seila Law} itself) acknowledges that independent agencies exercise executive power, \textit{Humphrey's Executor} ``does not even satisfy its own exception'' (\href{https://www.law.cornell.edu/supremecourt/text/19-7}{Cornell LII}).

On Part IV, however, Thomas dissented from the Court's severability holding. Drawing on traditional Article III limits and scholarship by Kevin Walsh and Jonathan Mitchell, Thomas argued that federal courts possess only the ``negative power to disregard an unconstitutional enactment'' in a particular case, not a legislative power to ``excise'' or rewrite statutory schemes based on speculative inquiries into hypothetical congressional intent (\href{https://www.law.cornell.edu/supremecourt/text/19-7}{Cornell LII}). Because Seila Law's constitutional injury persisted through Director Kraninger's continued litigation of the CID while statutorily insulated from removal, Acting Director Mulvaney's prior ratification could not cure the defect. Thomas would have skipped severability entirely and simply denied the CFPB's petition to enforce the CID (\href{https://www.law.cornell.edu/supremecourt/text/19-7}{Cornell LII}).

\subsection{Concurrence in the Judgment as to Severability and Dissent in Part --- Justice Kagan (Joined by Justices Ginsburg, Breyer, and Sotomayor)}

Justice Kagan delivered a vigorous functionalist and historical dissent on the merits, while concurring in the judgment on severability (giving Part IV's outcome a 7--2 majority) (\href{https://www.law.cornell.edu/supremecourt/text/19-7}{Cornell LII}). First, Kagan argued that the constitutional text says nothing about presidential removal; instead, the Necessary and Proper Clause (Art. I, \S 8, cl. 18) and the Appointments Clause grant Congress broad authority to structure executive departments and define officer tenure, so long as Congress does not impede the President's ability to perform his own constitutional duties (such as foreign affairs and military command). The Take Care Clause imposes a fiduciary duty to ensure faithful execution of the laws, not an illimitable power to fire for policy disagreements (\href{https://www.law.cornell.edu/supremecourt/text/19-7}{Cornell LII}).

Second, Kagan challenged the majority's history, contending that the Framers held divergent views (Hamilton in Federalist No. 77 assumed Senate consent was required for removal; Madison in Federalist No. 39 wrote that tenure of ministerial offices was a subject of legal regulation), that the Decision of 1789 was deeply divided and never settled Congress's power to regulate tenure, and that Congress from the Founding onward insulated financial regulators---from the Comptroller of the Treasury and the Second Bank of the United States to the Civil War Comptroller of the Currency, the ICC, the Federal Reserve, the FTC, and the SEC (\href{https://www.law.cornell.edu/supremecourt/text/19-7}{Cornell LII}).

Third, Kagan argued that \textit{Humphrey's Executor}, \textit{Wiener v. United States}, 357 U.S. 349 (1958), and \textit{Morrison} established a general rule permitting for-cause removal protections unless they ``impede the President's ability to perform his constitutional duty,'' not a rigid rule of at-will removal with two ``gerrymandered'' exceptions. Finally, she attacked the majority's single-Director vs. multimember-board distinction as backwards: a single Director is \textit{easier} for a President to oversee, persuade, and hold accountable than a ``gaggle'' of five staggered, bipartisan commissioners, and novelty is no constitutional defect under the Necessary and Proper Clause (\href{https://www.law.cornell.edu/supremecourt/text/19-7}{Cornell LII}).

\section{Examples: Future Applications}

\begin{enumerate}
\item \textbf{Same-Side (Unconstitutional): Single Director of the Federal Housing Finance Agency (FHFA).} Congress creates the FHFA to regulate Fannie Mae and Freddie Mac, led by a single Director serving a five-year term removable only ``for cause.'' Even though the FHFA regulates government-sponsored enterprises rather than purely private consumers, \textit{Seila Law} squarely controls: the Director is a single principal officer wielding significant executive power (including acting as conservator and executing financial agreements). Because the agency fits neither the \textit{Humphrey's Executor} multimember exception nor the \textit{Morrison} inferior-officer exception, the for-cause removal restriction violates Article II. (The Supreme Court reached precisely this result in \textit{Collins v. Yellen}, 594 U.S. 220 (2021)) (\href{https://www.law.cornell.edu/wex/unitary_executive_theory_\%28uet\%29}{Cornell LII Wex}).

\item \textbf{Same-Side (Unconstitutional): Independent Single-Head Environmental Enforcement Czar.} Suppose Congress restructures the Environmental Protection Agency (EPA) so that its single Administrator serves a six-year term removable only for ``inefficiency, neglect of duty, or malfeasance in office'' while retaining authority to promulgate binding nationwide emissions rules and seek civil penalties in federal court. Under \textit{Seila Law}, this structure is plainly unconstitutional: it concentrates massive executive rulemaking and enforcement authority in a single principal officer insulated from presidential control, with no historical pedigree and no multimember deliberative check.

\item \textbf{Opposite-Side (Constitutional under \textit{Seila Law}): Traditional Bipartisan Multimember Commission with Primarily Adjudicatory/Investigative Functions.} Suppose Congress creates a five-member Federal Maritime Claims Commission---with no more than three members from the same political party, serving staggered seven-year terms removable only for inefficiency, neglect of duty, or malfeasance---whose duties are limited to adjudicating statutory shipping-fee disputes and submitting reports to Congress, without bringing civil-penalty enforcement suits in federal court. Under \textit{Seila Law}'s express preservation of \textit{Humphrey's Executor} and \textit{Wiener}, this agency falls squarely within the multimember expert-body exception and survives Article II scrutiny.

\item \textbf{Opposite-Side (Constitutional under \textit{Seila Law}): Single Inferior Officer with Narrow, Supervised Duties.} Suppose Congress authorizes the Attorney General to appoint a Special Inspector for Pandemic Relief Fraud---an inferior officer removable by the Attorney General only for good cause---whose mandate is limited to auditing a single statutory grant program for three years and referring suspected fraud to the Department of Justice, without general rulemaking authority or power to set nationwide enforcement policy. Because the officer is an inferior officer with narrowly tailored duties and no broad policymaking or administrative authority, this structure falls within the \textit{Perkins}/\textit{Morrison} exception expressly left intact by \textit{Seila Law}.

\item \textbf{Fence-Sitter: Modern Multimember Independent Regulatory Agencies Wielding Substantial Executive Enforcement Power (e.g., the Modern FTC or SEC).} Consider a challenge to the for-cause removal protections of the five Commissioners of the modern Federal Trade Commission or the Securities and Exchange Commission, which---unlike the 1935 FTC as described in \textit{Humphrey's Executor}---now promulgate sweeping substantive rules and prosecute civil enforcement actions seeking billions of dollars in penalties in federal court. On one hand, \textit{Seila Law} explicitly refused to overrule \textit{Humphrey's Executor} and emphasized that multimember, partisan-balanced, staggered-term boards avoid concentrating power in one person. On the other hand, Chief Justice Roberts narrowed \textit{Humphrey's Executor} to multimember expert agencies ``that do not wield substantial executive power'' and noted in footnote 2 that the 1935 Court's premise that the FTC exercised no executive power ``has not withstood the test of time,'' while Justice Thomas's concurrence openly invited overruling \textit{Humphrey's Executor}. (Indeed, in \textit{Trump v. Slaughter} (2026), the Court ultimately confronted this exact fault line and invalidated the FTC's for-cause removal protection) (\href{https://www.law.cornell.edu/wex/unitary_executive_theory_\%28uet\%29}{Cornell LII Wex}; \href{https://www.law.nyu.edu/sites/default/files/2026-08/KKR\%2C-Officer-Removal-in-the-Early-Republic-20260731.pdf}{NYU Law}).
\end{enumerate}

\section{Critique}

\textbf{Originalist and Textualist Perspectives:} Originalist scholars are sharply divided over \textit{Seila Law}. Proponents of unitary executive theory, such as Aditya Bamzai and Saikrishna Bangalore Prakash, defend the Court's core premise, arguing from Founding-era texts, the Decision of 1789, and early presidential practice that Article II's ``executive Power'' encompassed an inherent presidential authority to remove executive officers at pleasure (\href{https://harvardlawreview.org/print/vol-136/the-executive-power-of-removal/}{Harvard Law Review}). From this formalist-originalist vantage point, \textit{Seila Law}'s chief shortcoming is that Chief Justice Roberts stopped halfway---preserving \textit{Humphrey's Executor} and severing the removal clause rather than adopting Justice Thomas's purer repudiation of independent agencies and modern severability doctrine. Conversely, a major wave of revisionist originalist and historical scholarship---including work by Jed Shugerman (\href{https://www.yalejreg.com/nc/a-reply-to-the-unitary-executive-theorists-on-the-misuse-of-historical-materials-by-jed-shugerman/}{Yale J. on Reg.}), Jane Manners and Lev Menand (\href{https://scholarship.law.columbia.edu/faculty_scholarship/3116/}{Columbia Law Scholarship Archive}), and Andrea Scoseria Katz and Noah Rosenblum (\href{https://harvardlawreview.org/forum/vol-136/removal-rehashed/}{Harvard Law Review Forum})---argues that the Court's history is flawed: the Decision of 1789 reached no consensus on constitutional indefeasibility, early Congresses routinely regulated officer tenure and vested removal or independent authority outside the President's plenary control, and ``inefficiency, neglect of duty, or malfeasance'' had specific historical meanings rooted in English and colonial officeholding rather than a blanket ban on presidential oversight.

\textbf{Progressive and Doctrinal Perspectives:} Progressive and functionalist scholars, such as Peter Shane (\href{https://www.acslaw.org/seila-law-lcc-v-cfpb-and-the-persistent-myths-of-presidentialism/}{ACS}) and administrative-law institutionalists, critique \textit{Seila Law} as judicial aggrandizement that substitutes formalist slogans for practical governance. They echo Justice Kagan's critique that the Necessary and Proper Clause empowers Congress to calibrate agency independence to protect technical financial regulation from short-term electoral manipulation and industry capture. Doctrinally, critics argue that the Court's bright line between single-head and multimember agencies is arbitrary and self-contradictory: if presidential accountability is the constitutional touchstone, a single Director is more transparent and responsive to presidential persuasion than a five-member bipartisan board with staggered terms. Critics also warn that elevating a rigid ``unitary executive'' rule threatens the stability of the modern administrative state, including financial stability regulators and civil-service protections.

\textbf{Underlying Competing Values:} At its core, \textit{Seila Law} pits \textit{democratic accountability and structural formalism} against \textit{institutional expertise, political insulation, and legislative flexibility}. The majority prioritizes a clear vertical chain of command running from every executive actor to the single nationally elected President, guarding against an unaccountable ``fourth branch'' exercising coercive state power. The dissent prioritizes Congress's pragmatic judgment in designing specialized, nonpartisan institutions insulated from White House interference, trusting inter-branch compromise over judicially enforced structural rigidity.

\section{Key Quotations}

\begin{quote}
``Under our Constitution, the `executive Power'---all of it---is `vested in a President,' who must `take Care that the Laws be faithfully executed.' Art. II, \S 1, cl. 1; id., \S 3.''
\end{quote}
\noindent --- Chief Justice Roberts, opinion of the Court (\href{https://www.law.cornell.edu/supremecourt/text/19-7}{Cornell LII}; \href{https://www.courtlistener.com/opinion/4764451/seila-law-llc-v-consumer-financial-protection-bureau/}{CourtListener}).

\begin{quote}
``The CFPB's single-Director structure contravenes this carefully calibrated system by vesting significant governmental power in the hands of a single individual accountable to no one.''
\end{quote}
\noindent --- Chief Justice Roberts, opinion of the Court (\href{https://www.law.cornell.edu/supremecourt/text/19-7}{Cornell LII}; \href{https://www.courtlistener.com/opinion/4764451/seila-law-llc-v-consumer-financial-protection-bureau/}{CourtListener}).

\begin{quote}
``We think it clear that Congress would prefer that we use a scalpel rather than a bulldozer in curing the constitutional defect we identify today.''
\end{quote}
\noindent --- Chief Justice Roberts, plurality opinion (Part IV) (\href{https://www.law.cornell.edu/supremecourt/text/19-7}{Cornell LII}; \href{https://www.courtlistener.com/opinion/4764451/seila-law-llc-v-consumer-financial-protection-bureau/}{CourtListener}).

\begin{quote}
``Continued reliance on \textit{Humphrey's Executor} to justify the existence of independent agencies creates a serious, ongoing threat to our Government's design. Leaving these unconstitutional agencies in place does not enhance this Court's legitimacy; it subverts political accountability and threatens individual liberty.''
\end{quote}
\noindent --- Justice Thomas, concurring in part and dissenting in part (\href{https://www.law.cornell.edu/supremecourt/text/19-7}{Cornell LII}; \href{https://www.courtlistener.com/opinion/4764451/seila-law-llc-v-consumer-financial-protection-bureau/}{CourtListener}).

\begin{quote}
``The majority's general rule does not exist. Its exceptions, likewise, are made up for the occasion---gerrymandered so the CFPB falls outside them.''
\end{quote}
\noindent --- Justice Kagan, concurring in the judgment with respect to severability and dissenting in part (\href{https://www.law.cornell.edu/supremecourt/text/19-7}{Cornell LII}; \href{https://www.courtlistener.com/opinion/4764451/seila-law-llc-v-consumer-financial-protection-bureau/}{CourtListener}).


\end{document}
